% !PRINTED 633
\weberpage{633}
These resolvents are integral algebraic numbers in certain
cyclotomic fields, and we have, when $\varepsilon^\lambda$ is not $=1$,
\[
(\varepsilon^\lambda,\varrho)\,(\varepsilon^{-\lambda},\varrho) = (-1)^\lambda p,
\]
and for $\lambda = m'n$
\weq{5}{(r^n,\eta)\,(r^{-n},\eta) = \pm\, p.}

All these resolvents are therefore divisors of the natural
prime $p$.

Besides these resolvents we have derived, at the place cited,
still other integers of the cyclotomic fields, namely
\weq{6}{\psi_{\lambda,\mu}(\varepsilon) = \sum\limits_{1,\,p-2}^{t} \varepsilon^{\mu\,\mathrm{ind}\,t - (\lambda+\mu)\,\mathrm{ind}\,(t+1)},}
(in which the indices refer to the base $g$), which in general belong to the field $\Omega_{p-1}$, and in the special case where $\lambda$ and $\mu$ are divisible by $m'$, to the field $\Omega_m$. Of particular importance among them are the numbers of the field $\Omega_m$:
\weq{7}{\psi_{sm',\,m'}(\varepsilon) = \psi_s(r) = \sum\limits_{1,\,p-2}^{t} r^{\mathrm{ind}\,t - (s+1)\,\mathrm{ind}\,(t+1)}.}

These numbers $\psi_{\lambda,\mu}(\varepsilon)$ too, when none of the numbers $\lambda$, $\mu$, $\lambda+\mu$ is divisible by $p-1$, are divisors of $p$ by the formula
\[
\psi_{\lambda,\mu}(\varepsilon)\,\psi_{\lambda,\mu}(\varepsilon^{-1}) = p,
\]
from which again, by (7), there results:
\weq{8}{\psi_s(r)\,\psi_s(r^{-1}) = p,}
provided that neither $s$ nor $s+1$ is divisible by $m$. Certain combinations of the resolvents (4), among them the $m$th power, we have expressed by means of the functions $\psi$ [Bd.\ I, \S.\,169, (14), (15)]:
\weq{9}{(r,\eta)^n\,(r^n,\eta)^{-1} = \psi_1(r)\,\psi_2(r)\cdots\psi_{n-1}(r),}
so long as $n<m$, and
\weq{10}{(r,\eta)\,(r^{-1},\eta) = \pm\,p}
\weq{11}{(r,\eta)^m = (-1)^{\frac{p-1}{m}}\, p\, \psi_1(r)\,\psi_2(r)\cdots\psi_{m-2}(r),}
from which it appears that these $m$th powers belong to the field $\Omega_m$.

We now denote, as before, by $\mathfrak{p}_n$ the various
prime divisors of $p$ in the field $\Omega_m$. The numbers $\psi_s(r)$, which likewise belong to this field and are factors of $p$, must
%CONT>>
